\scanpage{50}

% B01 | head, proof of implication and inserted hypothesis
$c\unsure{'|)}\Rightarrow b'$ Soit en effet $\Gamma$ un groupe comme dans $b')$, \inserted{Supposons d'abord $\ell\ne p$.} Soit $A=\bar k(\Gamma)$ l'algèbre de $\Gamma$ à coefficients dans la clôture algébrique $\bar k$ de $k$, \inserted{$\mathfrak g_k$ y opérant de façon évidente.} C'est une algèbre déf$/k$, commutative et \deleted{semi-simple} sép\uline{ara}ble (car $[\Gamma:e]$ est premier à $p$), donc en vertu de $c')$, on a $H^i(k,G)=0$ si $i>n$, $G$ désignant le groupe des éléments inversibles de $A$.

% B02 | upper middle, norm and its kernel
Soit $G_1$ le sous-groupe de $G$ formé des $x\in G$ tels que $N_{A/\bar k}x=1$ (où \deleted{dét. \unsure{$=$|est}} $N_{A/k}x=$ dét. de la translation à g. par $x$ dans $A$), c'est un sous-groupe \uline{défini sur $k$} (noyau de l'homomorphisme $N_{A/k}:G\longrightarrow\bar k^*$, qui est déf$/k$).

% B03 | middle, simple module and intersection
Comme $\Gamma$ est \uline{simple} comme $\mathfrak g$-module, \inserted{et que $\Gamma\subset G$} il en résulte que \deleted{\unsure{$N_{A/k}$}} l'on a $\Gamma\subset G_1$, ou $\Gamma\cap G_1=\deleted{\unsure{\{1\}|\{0\}}}\{e\}$, et on vérifie facilement que ce deuxième cas est exclu (distinguer le cas \deleted{\unsure{pair|$p$|premier}}\inserted{$\ell$} impair, et le cas $\ell=2$ $\deleted{\unsure{8|\}}}$, on note \deleted{\unsure{que}} dans ce dernier cas qu'en tous cas on a $N_{A/k}\xi=\pm1$ si $\xi\in\Gamma$).

% B04 | lower, order of the multiplication kernel
Donc $\Gamma\subset G_1$. Comme $\ell\Gamma\deleted{\unsure{G|=0}}=\{0\}$, il en résulte que $\Gamma$ est contenu dans le noyau de l'hom. $G_1\xrightarrow{\ell}G_1$. Ce noyau est d'ailleurs d'ordre $\ell^{n_1}$, où $n_1=n-1$ est la dim. de $G_1$, donc $n$ l'ordre de $G$.

% archive sheet number: 49

